% ---------------------------------------------------------------
% Section 4 — Implementation in the BX3 Framework
% Paper: The Bailout Protocol: Structural Compulsion for Exception Escalation in Autonomous Systems
% ---------------------------------------------------------------

\section{Implementation in the BX3 Framework}
\label{sec:bailout-implementation}

The Bailout Protocol is not a supplementary safety feature layered onto an existing autonomous system. Within the BX3 framework, it is the \emph{runtime expression} of the framework's upstream accountability guarantee: the point at which a declarative commitment to human oversight becomes an architecturally enforced, blocking, auditable operational decision. This section details how the protocol integrates with the Purpose Layer as the human accountability anchor (\S\ref{sec:purpose-layer-anchor}), how the Bounds Engine detects boundary transgression and triggers the bail-out cascade (\S\ref{sec:bounds-engine-triggering}), how the Fact Layer enforces a complete forensic audit trail (\S\ref{sec:fact-layer-enforcement}), how the forensic ledger records the full bail-out lifecycle (\S\ref{sec:forensic-ledger-recording}), and how the protocol's state machine functions as a natural extension of the Fact Layer's evidentiary architecture (\S\ref{sec:state-machine-extension}). We establish that the Bailout Protocol is the necessary and sufficient operationalisation of the upstream guarantee: without it, the framework's declarative accountability commitments remain unsatisfied at runtime.

% ----------------------------------------------------------------
\subsection{Purpose Layer: Human Accountability Anchor}
\label{sec:purpose-layer-anchor}

The BX3 framework enforces a non-negotiable architectural rule: the Purpose Layer must always remain human-anchored. This is not a convention or a preferred configuration — it is enforced at the Fact Layer's structural constraints. The Purpose Layer encodes the terminal objectives and normative boundaries of the system as structured, cryptographically signed declarations authored by designated human accountable parties. These declarations carry explicit moral and legal weight; they are not heuristics derived from system observation, and they cannot be overridden by any machine actor in the hierarchy under any operational circumstance.

The Purpose Layer provides three distinct and non-substitutable functions within the Bailout Protocol. First, it defines the \emph{terminal state conditions}: the precise circumstances under which continued autonomous operation would contradict the system's declared purpose. These conditions are grounded in explicit intent declarations, not in probabilistic thresholds — the distinction between a recommendation and a prohibition is made by human authority, not by inference. Second, the Purpose Layer establishes the \emph{accountability chain}: a structured mapping of which human parties are responsible for which classes of outcomes, and what obligations survive a bail-out event. When a bail-out occurs, no accountability vacuum exists — some designated human receives the forensic record and bears the obligation to act. Third, the Purpose Layer issues the \emph{signed capability grant} that authorises the Bailout Protocol to act on behalf of the accountable parties when terminal conditions are met.

This signed capability grant is the critical mechanism that prevents the Bailout Protocol from functioning as an autonomous override. The protocol does not act spontaneously; it acts on assignment, within the scope of an authorisation that originates from human authority. The grant is time-bounded, scope-limited, and revocable. If the Purpose Layer revokes the grant, the Bailout Protocol cannot activate — no autonomous escalation can occur in an accountability vacuum. This property is architecturally enforced: the Fact Layer will not record a bail-out transition without a valid, non-revoked capability grant from the Purpose Layer present in the chain of facts. The Purpose Layer thus provides the semantic grounding for every bail-out event: it answers not only ``who is responsible?'' but also ``under which governing commitment was this escalation mandatory?''

In operational terms, the Purpose Layer exposes three canonical fields to the Bailout Protocol. The \texttt{nodeId} field identifies the agent or process originating the escalation. The \texttt{humanRoot} field designates the human decision-maker at the root of the accountability chain — the terminal recipient of all unresolvable escalations. The \texttt{accountabilityAnchor} field provides a stable reference string that labels the governing policy under which this node operates, enabling post-hoc correlation between bail-out events and the governance instruments that authorised the agent's operation. When a child node is spawned, it inherits a \texttt{parentId} and a \texttt{parentPurposeHash} that cryptographically binds it to its parent's purpose declaration. The chain is tamper-evident: any divergence between a child's purpose and its parent's purpose is detectable by comparing the stored hash. The parent purpose hash is established at spawn time and is immutable for the node's operational lifetime — neither the child agent nor any subsequent actor in the system can modify or obscure the parent link after creation.

The Purpose Layer does not monitor the state machine, does not evaluate trigger conditions, and does not participate in the execution of resolution procedures. These are Bounds Engine and Fact Layer responsibilities. The Purpose Layer's role is strictly declarative: it defines what the system is for, who is accountable, and what authorisations it has granted. This separation of concerns is architecturally enforced. A machine actor that compromises the Purpose Layer would not thereby gain the ability to override the Bailout Protocol, because the Fact Layer independently verifies the existence and validity of the capability grant before recording any bail-out transition. The human anchor is structurally irreplaceable by any machine actor in the hierarchy.

% ----------------------------------------------------------------
\subsection{Bounds Engine: Bail-Out Triggering Authority}
\label{sec:bounds-engine-triggering}

The Bounds Engine is the active constraint layer of the BX3 framework. Where the Purpose Layer encodes what an agent is \emph{authorised} to do, the Bounds Engine encodes what an agent is \emph{not permitted} to do, what conditions constitute a violation, and what response is mandated when a violation is detected. The Bounds Engine is the triggering authority for the Bailout Protocol: it holds the encoded boundary conditions, monitors agent behaviour against those conditions in real time, and issues the bail-out signal when any boundary is crossed or when the agent encounters a condition it cannot resolve within its authorised scope.

The Bounds Engine operates on a three-stage model within the bail-out lifecycle. The first stage is \emph{bound encoding}: boundary conditions are expressed as executable predicates over the agent's observable state and action space. These predicates are authored during agent configuration and stored in the Bounds Registry. Example bounds include: ``do not authorise transactions exceeding \$10,000 without human confirmation,'' ``do not send communications outside the defined stakeholder list,'' ``halt if factual inconsistency exceeds a confidence threshold of 0.85,'' and ``do not spawn a child agent at depth $d \geq D_{\max}$.'' Each bound is assigned a severity classification that determines the urgency of the associated bail-out pathway.

The second stage is \emph{bound monitoring}: the Bounds Engine runs as a sidecar process that receives telemetry from the agent's action pipeline. Telemetry is evaluated against the encoded predicates at a configurable frequency — typically per-action for high-stakes boundaries and periodic for behavioural boundaries. The engine does not block agent actions synchronously; it evaluates them asynchronously to minimise latency impact on the agent's primary operation. When a predicate evaluates to \texttt{true}, indicating a boundary violation, the Bounds Engine transitions immediately to the third stage.

The third stage is \emph{bound enforcement}: the Bounds Engine emits a \texttt{BoundViolationEvent} and simultaneously dispatches the bail-out signal to the agent's execution context. The signal includes the bound identifier, the violated predicate, a Mandate reference, and a severity classification. Critically, the Bounds Engine does not attempt to resolve the violation — resolution is the purpose of the bounded resolution window in the Bailout Protocol state machine. The Bounds Engine's role is detection and signalling, not adjudication. This separation is architectural: if the Bounds Engine were permitted to both detect and resolve violations, the upstream accountability guarantee would be contingent on the correctness of the same actor that detected the violation.

The Bounds Engine also triggers the Bailout Protocol when an agent encounters a condition it cannot classify or resolve — even when no explicit bound has been violated. This scenario covers the case of \emph{unknown unknowns}: situations where the agent's purpose projection does not contain sufficient information to authorise a response. The agent submits a resolution request to the Bounds Engine; the Bounds Engine evaluates the request against the encoded bounds; if no bound applies but the agent has been unable to resolve the condition after exhausting its local resolution procedures, the Bounds Engine issues a \texttt{bailout\_signal} with reason \texttt{unresolvable\_condition}. This trigger condition ensures that the Bailout Protocol is activated not only by known-bound violations but also by the exhaustion of the agent's capability envelope, preventing the system from drifting into autonomous resolution of conditions outside its certified scope.

% ----------------------------------------------------------------
\subsection{Fact Layer: Audit Trail Enforcement}
\label{sec:fact-layer-enforcement}

The Fact Layer is the deterministic enforcement substrate of the BX3 framework. It is the layer that \emph{decides}, in the sense that no proposed action becomes real until the Fact Layer has evaluated it against pre-committed constraints and cleared it for execution. Within the Bailout Protocol, the Fact Layer performs three distinct enforcement functions: a pre-commit hook that validates bail-out transitions before they are recorded, a completeness invariant that ensures the audit chain is never broken, and a hard-block function that prevents any machine actor from bypassing or modifying the protocol's execution record.

The Fact Layer's pre-commit hook is the mechanism that makes the Bailout Protocol's state transitions architecturally compulsory rather than operationally optional. Before any state transition is recorded, the Fact Layer verifies that the transition satisfies three conditions: first, that a valid capability grant from the Purpose Layer exists and has not been revoked; second, that the transition is consistent with the transition relation defined by the protocol's state machine; and third, that the transition's triggering event is present in the Bounds Engine's registry and has not been superseded by a later resolution. If any of these conditions fails, the Fact Layer rejects the transition, logs a \texttt{protocol-violation} entry, and surfaces an alert to the human accountability anchor designated in the Purpose Layer.

The Fact Layer also enforces a \emph{completeness invariant}: for every Bailout Protocol event, a corresponding chain of facts must exist from the triggering boundary event through to the terminal state resolution. If this chain is broken — whether due to data corruption, system crash, or deliberate excision — the Fact Layer raises a structural integrity alert. This alert is surfaced to the human auditors designated in the Purpose Layer's accountability chain. The alert itself is logged as a fact, ensuring that even a compromise of the audit infrastructure is itself auditable.

The Fact Layer does not permit retroactive modification of any Bailout Protocol record. Once a \texttt{BailoutLedgerEntry} is committed, it is append-only. No actor — including the Purpose Layer post-provisioning, the Bounds Engine, any administrative process, or any external principal — may execute an operation that overwrites, redirects, or deletes a committed entry. The immutability is protocol-enforced, not access-control-enforced. Unlike access-control models where immutability can be circumvented by privilege escalation, the Fact Layer write protocol makes modification structurally impossible regardless of credential compromise or software vulnerability.

% ----------------------------------------------------------------
\subsection{Forensic Ledger Recording}
\label{sec:forensic-ledger-recording}

The Forensic Ledger is the append-only record of all Bailout Protocol events. It is implemented as a designated partition within the Fact Layer's hash-chained record store, optimised for bail-out event retrieval and analysis. Every state transition, trigger evaluation, timeout event, pathway selection decision, parent invocation, human acknowledgement, and human directive is committed to the Forensic Ledger as an append-only entry \emph{before} any subsequent action is taken. Entries are chained by SHA-256 hash: each entry $f$ contains $\text{hash}(f-1)$, creating a tamper-evident structure that makes retrospective modification of any entry computationally detectable by any observer with ledger read access.

The minimum entry types generated by a single Bailout Protocol execution are:

\begin{itemize}\itemsep=0pt
\item[\texttt{bailout-trigger}] Originating node identifier, exception identifier, trigger condition sub-components, confidence score.
\item[\texttt{bounds-resolution-attempt}] Each resolution procedure submitted during the bounded resolution window, with authorisation ledger reference and Fact Layer approval or rejection outcome.
\item[\texttt{bailout-signal}] Explicit \texttt{bailout\_signal} emissions with Bounds Engine's stated reason.
\item[\texttt{timeout-bounds-expired}] Expiration of the bounded resolution timeout $T_{\text{bounds}}$ without resolution.
\item[\texttt{parent-escalation}] Each \texttt{SendToParent} invocation, including parent node identifier and delivery confirmation.
\item[\texttt{human-acknowledged}] Arrival at \texttt{humanRoot}, pending human decision.
\item[\texttt{human-directive}] Human directive text, issuing identity, and timestamp.
\item[\texttt{resolved}] Terminal state with full end-to-end latency from \texttt{firedAt} to \texttt{resolvedAt}.
\end{itemize}

Each \texttt{BailoutLedgerEntry} is a composite record that consolidates the full bail-out context into a single queryable entity. The entry contains: an entry identifier (monotonically increasing, enabling strict total ordering across all agents in the system); the agent identity with its active Purpose Mandate reference; the full bound violation record including predicate, threshold, and observed value; the complete propagation chain as a sequence of node identifiers; and the terminal state resolution with latency measurement. The combination of the \texttt{propagationChain} array inside the \texttt{BailoutTrigger} record and the corresponding \texttt{LedgerEntry} sequence in the \texttt{ForensicLedger} provides dual-mode auditability: the chain gives a human-readable provenance graph, while the ledger gives a machine-verifiable cryptographic proof of integrity.

An auditor who queries the Forensic Ledger for any trigger whose \texttt{triggerCondition} is \texttt{accountability\_boundary} can confirm that the propagation chain terminates at the declared \texttt{humanRoot} with a corresponding \texttt{resolved} entry. The completeness of this chain — from trigger to human acknowledgement to resolution — is what distinguishes the Bailout Protocol's audit trail from conventional log files, which record events but do not enforce the structural relationship between them.

% ----------------------------------------------------------------
\subsection{State Machine as Fact Layer Extension}
\label{sec:state-machine-extension}

The Bailout Protocol's state machine is not embedded within the agent's execution context, where it could be subject to the same faults or compromised conditions that necessitate the bail-out. Instead, it is implemented as a fact-driven observer operating as a natural extension of the Fact Layer: the Bounds Engine and the Forensic Ledger operate as independent processes that observe the agent's telemetry and record facts about its state transitions. The agent's own execution context is unaware of the state machine; it simply receives and honours bail-out signals. This architectural separation ensures that the bail-out mechanism cannot be circumvented by a compromised or malfunctioning agent.

The state machine comprises five states. \texttt{AUTHORIZED} is the default operational state: the agent is operating within its active Purpose Mandate, all encoded bounds are satisfied, and no bail-out condition is present. \texttt{BOUND\_VIOLATION\_DETECTED} is a transitional state that captures the window between detection of a boundary violation and receipt of the bail-out signal: the Bounds Engine has emitted a \texttt{BoundViolationEvent} but the agent has not yet acknowledged the bail-out. \texttt{ESCALATING} is entered when the Bounds Engine issues a \texttt{bailout\_signal} after the bounded resolution window expires without resolution: the agent has received the signal and entered a halted state, and the Fact Layer pre-commit hook has verified the transition. \texttt{HUMAN\_ACKNOWLEDGED} is entered when the escalation reaches the declared \texttt{humanRoot} and the human accountable party has acknowledged receipt: no autonomous resolution is permitted beyond this state. \texttt{RESOLVED} is the terminal state: the human has issued a directive and the Fact Layer has recorded the resolution.

Critically, state transitions are themselves recorded as facts in the Fact Layer. Each transition produces an \texttt{ImmutableFact} with the transition type, the source and destination states, the triggering event identifier, and a timestamp. This means the state machine's own behaviour is subject to the same immutability and cryptographic chaining as the agent's operational facts. There is no privileged execution path that bypasses the fact chain — the state machine's integrity is itself part of what the audit trail verifies. The Bailout Monitor, a designatedFact Layer process, ensures that the protocol selects functional pathways even under partial failure conditions: degraded or unresponsive pathways are excluded from selection until a health check succeeds. The Fact Layer's embedded position of the state machine ensures that the Bailout Protocol's guarantees are not contingent on the correct behaviour of any machine actor.

% ----------------------------------------------------------------
\subsection{Bailout Protocol as Runtime Expression of the Upstream Accountability Guarantee}
\label{sec:upstream-guarantee}

BX3 makes an explicit upstream guarantee: that the system's behaviour will remain within the bounds declared by its human accountable parties, and that any transgression beyond those bounds will be detected, recorded, and resolved in a manner that preserves accountability. The Bailout Protocol is the runtime expression of this guarantee. It is the mechanism by which an abstract, declarative accountability promise becomes an enforceable, observable property of a running system.

The relationship can be formalised. Let $P$ represent the Purpose Layer's declarative constraints. Let $B$ represent the Bounds Engine's boundary detection and simulation logic. Let $F$ represent the Fact Layer's forensic ledger. Let $S$ represent the current system state. BX3's upstream guarantee states:

\begin{equation}
    \forall S,\; \text{if}\; S \models \neg P\; \text{then}\;
    \exists e \in B : e.\text{trigger} \Rightarrow
    \exists f \in F : f.\text{records}(e) \land \text{BailoutProtocol}(e).\text{resolves}()
\label{eq:upstream-guarantee}
\end{equation}

In plain terms: for any system state that violates the Purpose Layer's constraints, the Bounds Engine will emit a triggering event, the Fact Layer will record that event in a complete chain, and the Bailout Protocol will resolve it. The guarantee is compositional: each layer depends on the others, and the failure of any layer to perform its designated function constitutes a breach of the upstream guarantee itself, surfacing accountability at the level of the BX3 framework rather than at the individual node.

The causal chain from declaration to enforcement operates as follows. The \texttt{humanRoot} field in the Purpose Layer \emph{declares} that a specific human is the final accountable authority for actions originating under this node's governance. The \texttt{accountabilityBoundary} condition in the Bounds Engine \emph{detects} when a proposed action falls outside the scope of autonomous authorisation. The Bailout Protocol's \texttt{fire()} method \emph{operationalises} the detection by creating a trigger that is obligated to propagate to the declared human. The state machine \emph{enforces} that no autonomous resolution is permitted beyond \texttt{HUMAN\_ACKNOWLEDGED}, making the guarantee non-negotiable. The Forensic Ledger \emph{records} every step of this chain, providing the evidence that the guarantee was honoured.

Without the Bailout Protocol, the upstream accountability guarantee would be a structural property with no enforcement mechanism. The human anchor $h(n)$ would be the designated terminus in principle, but no architectural compulsion would require exceptions to reach it. Machine actors could suppress, redirect, or silently absorb exception states, and the guarantee would hold only for systems whose machine actors chose to comply voluntarily. The upstream accountability guarantee would be a declaration, not a mechanism.

The Bailout Protocol converts the guarantee from a structural aspiration into an architecturally enforced compulsion through three coordinated properties. First, the deterministic transition relation ensures that escalation is not a routing decision made by a machine actor — it is a state machine transition that fires compulsorily when a trigger condition is satisfied and the bounds timeout expires. Second, the Fact Layer's pre-commit hook ensures that no machine actor can submit a retroactive resolution, extend a timeout, or redirect the propagation pathway. Third, the immutable parent-pointer architecture ensures that the propagation substrate itself cannot be rewired by any machine actor — the parent chain is established at spawn time and is read-only for the node's lifetime.

The Forensic Ledger entries generated by the protocol provide the retrospective proof of this guarantee. An auditor who accepts the BX3 Purpose Layer's formal guarantee as a design commitment can verify its enforcement by querying the Forensic Ledger for any trigger whose \texttt{triggerCondition} is \texttt{accountability\_boundary} and confirming that the \texttt{propagationChain} terminates at the declared \texttt{humanRoot} with a corresponding \texttt{RESOLVED} entry. The Bailout Protocol thus transforms an architectural commitment into a measurable, auditable property of every BX3 node — provable, not merely believed.